\scanpage{155}
UNIVERSITÉ DES SCIENCES ET TECHNIQUES\\
DU LANGUEDOC\\
----

Service du Personnel\\
Diffusion liste D + affichage\\
N° 4782

\uline{ACCES DES VEHICULES AUTOMOBILES}\\
\uline{SUR LE CAMPUS DE L'U.S.T.L.}\\
--------

Je vous rappelle que l'accès des véhicules automobiles sur le campus est réservé aux détenteurs d'un coupe-file.

Ceux-ci sont délivrés par le service du Personnel (bureau d'accueil) :

-- soit à titre définitif (macarons plastifiés auto-collants) pour les personnels permanents d'Enseignement, Administratifs, Techniques et de Service.

-- soit sous forme de laissez-passer provisoires (bristol vert) valables pour l'année universitaire. Ces derniers peuvent être délivrés :

\[
\left.\begin{array}{l}
\text{-- aux personnels du C.N.R.S.}\\
\text{-- aux personnels sous convention}
\end{array}\right\}
\begin{array}{l}
\text{au vu d'une attestation du Chef de}\\
\text{Service et de la carte grise du}\\
\text{véhicule}
\end{array}
\]

-- aux thésards d'Etat et de 3ème cycle, au vu d'un certificat de scolarité et de la carte grise du véhicule.

La validité des laissez-passer provisoires délivrés au titre de l'année universitaire 1979 - 1980 expirant le \uline{31 Octobre 1980,} les détenteurs de ces cartes sont invités à se présenter au Service du Personnel (bureau d'accueil) munis des pièces mentionnées ci-dessus, pour retirer les nouveaux laissez-passer valables pour l'année universitaire 1980 -- 1981.

Le Président,

L. Thaler\\
L. THALER

\figure{\begin{tikzpicture}[x=1cm,y=1cm]
\draw (15:1.55) arc[start angle=15,end angle=-298,radius=1.55];
\draw (15:1.16) arc[start angle=15,end angle=-291,radius=1.16];
\node[rotate=48] at (-1.05,.95) {\unsure{PRÉSIDENCE}};
\node[rotate=42] at (.65,-1.12) {\unsure{MONTPELLIER}};
\node[rotate=-22,align=center,font=\small] at (0,0) {\unsure{UNIVERSITÉ}\\\unsure{DES}\\SCIENCES et \unsure{TECHNIQUES}\\\unsure{DU}\\LANGUEDOC};
\end{tikzpicture}}

\margin{155}
